\chapter{Đặc tả yêu cầu phần mềm}

\section{Mục tiêu của chương}

Chương này trình bày quá trình khảo sát nghiệp vụ, xác định stakeholder, chốt yêu cầu chức năng, yêu cầu phi chức năng, phạm vi hệ thống và ma trận truy vết yêu cầu. Mục tiêu của chương không chỉ là liệt kê chức năng, mà còn chứng minh các chức năng trong hệ thống xuất phát từ nhu cầu vận hành của phòng mạch tư nhân.

\section{Bối cảnh và vấn đề của khách hàng}

Trong phạm vi đồ án môn học, khách hàng được mô phỏng là một phòng mạch tư nhân quy mô nhỏ đến trung bình. Quy trình vận hành điển hình của phòng mạch gồm: bệnh nhân đến đăng ký khám, lễ tân tạo lượt khám và số thứ tự, bác sĩ tiếp nhận khám, lập phiếu khám và kê đơn, thu ngân lập hóa đơn và ghi nhận thanh toán, quản lý theo dõi báo cáo hoạt động.

Các vấn đề thường gặp nếu quản lý thủ công gồm:

\begin{itemize}
    \item Hồ sơ bệnh nhân được lưu bằng sổ sách hoặc file rời rạc, khó tra cứu và dễ trùng lặp.
    \item Không có hệ thống cấp số thứ tự tự động, dễ gây nhầm lẫn trong quá trình tiếp nhận.
    \item Phiếu khám và đơn thuốc viết tay gây khó đọc, khó lưu trữ và khó tra cứu lịch sử.
    \item Hóa đơn tính thủ công dễ sai sót khi tổng hợp phí khám, tiền thuốc và các dịch vụ liên quan.
    \item Quản lý khó theo dõi lượt khám, doanh thu và tình hình sử dụng thuốc/dịch vụ.
\end{itemize}

\section{Khảo sát khách hàng giả định}

Do đây là đồ án môn học, phần khảo sát khách hàng được mô phỏng dựa trên nghiệp vụ phòng mạch tư nhân. Không có phỏng vấn khách hàng thật. Nhóm phân tích nhu cầu của các vai trò thường gặp trong phòng mạch và chuyển các nhu cầu đó thành yêu cầu phần mềm.

\begin{longtable}{|p{3.2cm}|p{4.3cm}|p{4.2cm}|p{4.2cm}|}
\caption{Tổng hợp stakeholder và nhu cầu nghiệp vụ}\\
\hline
\textbf{Đối tượng} & \textbf{Nhu cầu chính} & \textbf{Vấn đề hiện tại} & \textbf{Yêu cầu mong muốn} \\
\hline
Chủ phòng khám / Quản lý & Theo dõi hoạt động và doanh thu & Khó tổng hợp doanh thu, lượt khám bằng thủ công & Báo cáo tháng, thống kê lượt khám, doanh thu \\
\hline
Bác sĩ & Tra cứu lịch sử bệnh nhân, khám và kê đơn & Khó tra cứu hồ sơ cũ, đơn thuốc viết tay & Xem lịch sử, lập phiếu khám, kê đơn điện tử \\
\hline
Lễ tân & Tiếp nhận nhanh, xếp hàng trật tự & Dễ nhầm thông tin, không có số thứ tự hệ thống & Quản lý hồ sơ, tạo lượt khám, cấp số thứ tự \\
\hline
Thu ngân & Tính hóa đơn đúng, ghi nhận thanh toán & Tính thủ công dễ sai sót & Hóa đơn điện tử, thanh toán nhiều lần \\
\hline
Admin & Quản lý tài khoản, phân quyền, cấu hình & Chưa kiểm soát rõ ai được dùng chức năng nào & RBAC, quản lý tài khoản, quy định \\
\hline
Y tá / Điều dưỡng & Đo sinh hiệu, hỗ trợ bác sĩ & Ghi nhận rời rạc hoặc bằng giấy & Lưu sinh hiệu, tính BMI tự động \\
\hline
Kỹ thuật viên xét nghiệm & Nhận mẫu và nhập kết quả & Quy trình xét nghiệm rời rạc & Worklist xét nghiệm, nhập và xác nhận kết quả \\
\hline
Dược sĩ & Cấp phát thuốc, quản lý tồn kho & Khó kiểm soát hạn dùng và số lượng tồn & Cấp phát theo đơn, quản lý lô thuốc, FEFO \\
\hline
\end{longtable}

\section{Phạm vi hệ thống}

\subsection{Trong phạm vi}

Hệ thống được xác định là ứng dụng web nội bộ, phục vụ người dùng trong phòng mạch. Phạm vi chính gồm 30 yêu cầu chức năng, chia thành nhóm Phase 1 và Phase 2. Các yêu cầu này được ánh xạ sang các use case tổng quan trong sơ đồ use case.

\begin{itemize}
    \item Phase 1: đăng nhập, phân quyền, tài khoản, bệnh nhân, lượt khám, phiếu khám, kê đơn, hóa đơn, thanh toán, danh mục, quy định, báo cáo tháng.
    \item Phase 2: lịch hẹn, hàng đợi, check-in, sinh hiệu, dịch vụ, xét nghiệm, kho thuốc, cấp phát, tổ chức phòng khám và audit log.
    \item Hệ thống hỗ trợ 8 vai trò: ADMIN, RECEPTIONIST, DOCTOR, CASHIER, MANAGER, NURSE, LAB\_TECHNICIAN, PHARMACIST.
\end{itemize}

\subsection{Ngoài phạm vi}

Các chức năng sau không thuộc phạm vi phiên bản hiện tại: cổng bệnh nhân tự tra cứu online, đặt lịch trực tuyến từ phía bệnh nhân, nhắc lịch qua SMS/email, tích hợp bảo hiểm, multi-branch/multi-tenant, telemedicine, mobile app, advanced analytics, Docker, CI/CD và production deployment.

\section{Yêu cầu chức năng}

\begin{longtable}{|p{1.6cm}|p{5.6cm}|p{3cm}|p{1.8cm}|p{1.6cm}|}
\caption{Danh sách yêu cầu chức năng REQ-01 đến REQ-30}\\
\hline
\textbf{ID} & \textbf{Yêu cầu} & \textbf{Actor} & \textbf{Ưu tiên} & \textbf{Phase} \\
\hline
REQ-01 & Đăng nhập bằng email/password & Tất cả roles & Cao & P1 \\
\hline
REQ-02 & Phân quyền theo vai trò & ADMIN & Cao & P1 \\
\hline
REQ-03 & Quản lý tài khoản, CRUD và khóa tài khoản & ADMIN & Cao & P1 \\
\hline
REQ-04 & Tìm kiếm bệnh nhân & RECEPTIONIST, DOCTOR & Cao & P1 \\
\hline
REQ-05 & Tạo hồ sơ bệnh nhân & RECEPTIONIST & Cao & P1 \\
\hline
REQ-06 & Tiếp nhận bệnh nhân, tạo visit & RECEPTIONIST & Cao & P1 \\
\hline
REQ-07 & Cấp số thứ tự tự động, không trùng & RECEPTIONIST & Cao & P1 \\
\hline
REQ-08 & Xem danh sách lượt khám theo ngày/trạng thái & DOCTOR, RECEPTIONIST & Cao & P1 \\
\hline
REQ-09 & Bác sĩ mở lượt khám & DOCTOR & Cao & P1 \\
\hline
REQ-10 & Lập phiếu khám, ghi triệu chứng và chẩn đoán & DOCTOR & Cao & P1 \\
\hline
REQ-11 & Xem lịch sử khám bệnh nhân & DOCTOR & Cao & P1 \\
\hline
REQ-12 & Kê đơn thuốc điện tử & DOCTOR & Cao & P1 \\
\hline
REQ-13 & Hoàn tất phiếu khám & DOCTOR & Cao & P1 \\
\hline
REQ-14 & Lập hóa đơn & CASHIER & Cao & P1 \\
\hline
REQ-15 & Ghi nhận thanh toán & CASHIER & Cao & P1 \\
\hline
REQ-16 & Tra cứu hóa đơn & CASHIER, MANAGER & Cao & P1 \\
\hline
REQ-17 & Thay đổi quy định phòng khám & ADMIN, MANAGER & Trung bình & P1 \\
\hline
REQ-18 & Quản lý danh mục bệnh & ADMIN, MANAGER & Trung bình & P1 \\
\hline
REQ-19 & Quản lý danh mục thuốc & ADMIN, MANAGER & Trung bình & P1 \\
\hline
REQ-20 & Báo cáo tháng cơ bản & MANAGER & Trung bình & P1 \\
\hline
REQ-21 & Đặt và quản lý lịch hẹn & RECEPTIONIST & Trung bình & P2 \\
\hline
REQ-22 & Hàng đợi điện tử & RECEPTIONIST, NURSE & Trung bình & P2 \\
\hline
REQ-23 & Check-in lịch hẹn tạo lượt khám & RECEPTIONIST & Trung bình & P2 \\
\hline
REQ-24 & Đo và lưu sinh hiệu, BMI tự tính & NURSE & Thấp & P2 \\
\hline
REQ-25 & Quản lý danh mục dịch vụ và chỉ định & DOCTOR, ADMIN & Thấp & P2 \\
\hline
REQ-26 & Quy trình xét nghiệm order, sample, result, verify & LAB\_TECHNICIAN, DOCTOR & Thấp & P2 \\
\hline
REQ-27 & Quản lý lô thuốc và tồn kho & PHARMACIST & Thấp & P2 \\
\hline
REQ-28 & Cấp phát thuốc theo FEFO & PHARMACIST & Thấp & P2 \\
\hline
REQ-29 & Quản lý khoa phòng và bác sĩ & ADMIN, MANAGER & Thấp & P2 \\
\hline
REQ-30 & Nhật ký hệ thống audit log & ADMIN & Thấp & P2 \\
\hline
\end{longtable}

\section{Yêu cầu phi chức năng}

\begin{table}[H]
\centering
\caption{Yêu cầu phi chức năng}
\begin{tabularx}{\textwidth}{|p{3.7cm}|X|}
\hline
\textbf{Nhóm yêu cầu} & \textbf{Mô tả} \\
\hline
Bảo mật & Hệ thống sử dụng JWT và RBAC để kiểm soát truy cập theo vai trò; mật khẩu được hash bằng bcrypt; access token hết hạn theo cấu hình \code{JWT\_ACCESS\_EXPIRES}, refresh token hết hạn theo \code{JWT\_REFRESH\_EXPIRES} và bị revoke khi logout. \\
\hline
Toàn vẹn dữ liệu & Các ràng buộc unique, foreign key và transaction được sử dụng để đảm bảo dữ liệu nhất quán; các nghiệp vụ tạo lượt khám, thanh toán, kích hoạt quy định và cấp phát thuốc phải cập nhật dữ liệu theo transaction. \\
\hline
Hiệu năng & Trong môi trường local demo, các API nghiệp vụ phổ biến như tra cứu bệnh nhân, danh sách lượt khám và chi tiết hóa đơn nên phản hồi dưới 500ms với dữ liệu seed; hệ thống hướng đến phục vụ tối thiểu 10 người dùng nội bộ đồng thời trong phạm vi phòng mạch nhỏ. \\
\hline
Dễ sử dụng & Giao diện được tổ chức theo vai trò, có sidebar, form validation, thông báo lỗi và success toast; các màn hình nghiệp vụ phải có trạng thái loading, empty và error rõ ràng. \\
\hline
Bảo trì & Backend tổ chức theo module, controller, service, DTO; frontend tổ chức theo feature; các business rules quan trọng đặt ở service layer để dễ kiểm thử và tránh phụ thuộc giao diện. \\
\hline
Mở rộng & Modular monolith cho phép bổ sung Phase 2 mà không thay đổi toàn bộ kiến trúc ban đầu; các nghiệp vụ mở rộng như lab, queue, pharmacy được tách module riêng. \\
\hline
Triển khai & Phiên bản hiện tại hỗ trợ local development với PostgreSQL, NestJS và Vite; Docker, CI/CD và production deployment là hướng phát triển sau, không thuộc phạm vi khẳng định của báo cáo. \\
\hline
\end{tabularx}
\end{table}

\section{Use case tổng quan}

\LandscapeFigure{figures/ch01-requirements/usecase-phase1.png}{Use case diagram Phase 1 -- nghiệp vụ lõi}{fig:usecase-phase1}
\LandscapeFigure{figures/ch01-requirements/usecase-phase2.png}{Use case diagram Phase 2 -- nghiệp vụ mở rộng}{fig:usecase-phase2}

\section{Mô tả use case trọng tâm}

\subsection{UC07 -- Tạo lượt khám và cấp số thứ tự}

\begin{table}[H]
\centering
\caption{Mô tả use case UC07 -- Tạo lượt khám và cấp số thứ tự}
\begin{tabularx}{\textwidth}{|p{3.5cm}|X|}
\hline
\textbf{Use case ID} & UC07 \\
\hline
\textbf{Actor} & RECEPTIONIST là actor chính; ADMIN có thể hỗ trợ trong trường hợp vận hành đặc biệt. \\
\hline
\textbf{Tiền điều kiện} & Người dùng đã đăng nhập; bệnh nhân đã có hồ sơ trong hệ thống; ngày khám nằm trong phạm vi tiếp nhận. \\
\hline
\textbf{Hậu điều kiện} & Visit được tạo, có \code{queueNumber} duy nhất theo ngày và có trạng thái ban đầu phù hợp với quy trình tiếp nhận. \\
\hline
\textbf{Luồng chính} & Lễ tân tìm hoặc chọn bệnh nhân; hệ thống hiển thị form tạo lượt khám; lễ tân xác nhận; backend kiểm tra bệnh nhân tồn tại, kiểm tra chưa có visit trùng trong ngày, kiểm tra quota theo quy định, sinh số thứ tự trong transaction và trả về thông tin visit mới. \\
\hline
\textbf{Luồng ngoại lệ} & Nếu bệnh nhân đã có lượt khám trong ngày, hệ thống trả 409 Conflict; nếu vượt quota ngày, hệ thống trả 409 Conflict; nếu bệnh nhân không tồn tại, hệ thống trả 404 Not Found. \\
\hline
\textbf{Yêu cầu liên quan} & REQ-04, REQ-05, REQ-06, REQ-07, REQ-08, BR-05, BR-06, BR-07. \\
\hline
\end{tabularx}
\end{table}

\subsection{UC13 -- Hoàn tất phiếu khám}

\begin{table}[H]
\centering
\caption{Mô tả use case UC13 -- Hoàn tất phiếu khám}
\begin{tabularx}{\textwidth}{|p{3.5cm}|X|}
\hline
\textbf{Use case ID} & UC13 \\
\hline
\textbf{Actor} & DOCTOR là actor chính. \\
\hline
\textbf{Tiền điều kiện} & Bác sĩ đã mở phiên khám; phiếu khám đang ở trạng thái có thể cập nhật; thông tin chẩn đoán tối thiểu đã được nhập. \\
\hline
\textbf{Hậu điều kiện} & Examination chuyển sang trạng thái hoàn tất; Visit tương ứng được cập nhật sang trạng thái hoàn tất để thu ngân có thể lập hóa đơn. \\
\hline
\textbf{Luồng chính} & Bác sĩ nhập triệu chứng, chẩn đoán, bệnh tham chiếu và đơn thuốc nếu có; bác sĩ chọn hoàn tất; backend kiểm tra trạng thái phiếu khám, kiểm tra dữ liệu bắt buộc, lưu thay đổi và cập nhật trạng thái visit trong transaction. \\
\hline
\textbf{Luồng ngoại lệ} & Nếu thiếu chẩn đoán, hệ thống trả 400 BadRequest; nếu phiếu khám đã hoàn tất hoặc bị hủy, hệ thống không cho cập nhật; nếu còn dịch vụ/xét nghiệm bắt buộc chưa hoàn tất, đây là rủi ro đã được ghi nhận để cải thiện. \\
\hline
\textbf{Yêu cầu liên quan} & REQ-09, REQ-10, REQ-12, REQ-13, BR-10, BR-12, BR-14. \\
\hline
\end{tabularx}
\end{table}

\subsection{UC15 -- Ghi nhận thanh toán}

\begin{table}[H]
\centering
\caption{Mô tả use case UC15 -- Ghi nhận thanh toán hóa đơn}
\begin{tabularx}{\textwidth}{|p{3.5cm}|X|}
\hline
\textbf{Use case ID} & UC15 \\
\hline
\textbf{Actor} & CASHIER là actor chính; MANAGER có quyền xem lại hóa đơn và báo cáo. \\
\hline
\textbf{Tiền điều kiện} & Lượt khám đã hoàn tất; hóa đơn đã được lập; hóa đơn chưa bị hủy và chưa thanh toán đủ. \\
\hline
\textbf{Hậu điều kiện} & Payment được ghi nhận; \code{paidAmount} và \code{InvoiceStatus} được cập nhật nhất quán. \\
\hline
\textbf{Luồng chính} & Thu ngân mở chi tiết hóa đơn, nhập số tiền và phương thức thanh toán; backend kiểm tra số tiền lớn hơn 0, không vượt phần còn lại, tạo payment và cập nhật trạng thái hóa đơn trong transaction. \\
\hline
\textbf{Luồng ngoại lệ} & Nếu số tiền không hợp lệ hoặc vượt phần còn lại, hệ thống trả 400 BadRequest; nếu hóa đơn đã PAID/VOID, hệ thống từ chối ghi nhận thanh toán mới. \\
\hline
\textbf{Yêu cầu liên quan} & REQ-14, REQ-15, REQ-16, BR-15, BR-17. \\
\hline
\end{tabularx}
\end{table}

\section{Acceptance criteria chọn lọc}

\begin{longtable}{|p{2cm}|p{12.8cm}|}
\caption{Acceptance criteria cho các yêu cầu quan trọng}\\
\hline
\textbf{REQ} & \textbf{Điều kiện chấp nhận} \\
\hline
REQ-01 & Đăng nhập đúng email/password trả về access token, refresh token và thông tin người dùng; đăng nhập sai trả 401 Unauthorized. \\
\hline
REQ-02 & Người dùng không đủ vai trò khi gọi API được bảo vệ phải nhận 403 Forbidden hoặc bị chặn bởi guard tương ứng. \\
\hline
REQ-03 & ADMIN có thể tạo, cập nhật, khóa tài khoản; tài khoản bị khóa không thể đăng nhập. \\
\hline
REQ-04 & Người dùng được phép có thể tìm bệnh nhân theo mã, tên, số điện thoại hoặc định danh; kết quả không hiển thị dữ liệu nhạy cảm không cần thiết. \\
\hline
REQ-05 & Tạo hồ sơ bệnh nhân thành công sinh \code{patientCode}; citizenId trùng bị từ chối bằng 409 Conflict. \\
\hline
REQ-06 & Một bệnh nhân không được có hai lượt khám khác nhau trong cùng một ngày. \\
\hline
REQ-07 & Hai lượt khám cùng ngày không được có cùng số thứ tự; số thứ tự được tạo trong transaction. \\
\hline
REQ-08 & Danh sách lượt khám lọc được theo ngày và trạng thái; bác sĩ và lễ tân nhìn thấy đúng phạm vi công việc. \\
\hline
REQ-09 & Bác sĩ chỉ mở phiên khám từ lượt khám hợp lệ; mở thành công tạo examination và cập nhật trạng thái visit. \\
\hline
REQ-10 & Phiếu khám lưu được triệu chứng, chẩn đoán, bệnh tham chiếu và phí khám; phiếu đã hoàn tất không được cập nhật tùy tiện. \\
\hline
REQ-12 & Đơn thuốc phải có thuốc active, số lượng hợp lệ và hướng dẫn dùng thuốc; không kê trùng cùng một thuốc trong một đơn. \\
\hline
REQ-13 & Không hoàn tất phiếu khám nếu chưa có chẩn đoán. \\
\hline
REQ-14 & Chỉ lập hóa đơn khi lượt khám đã hoàn tất; hóa đơn phản ánh phí khám, thuốc và dịch vụ liên quan. \\
\hline
REQ-15 & Số tiền thanh toán không được vượt quá số tiền còn lại của hóa đơn. \\
\hline
REQ-16 & Thu ngân và quản lý tra cứu được hóa đơn theo trạng thái/ngày; thông tin thanh toán được cập nhật sau mỗi payment. \\
\hline
REQ-28 & Cấp phát thuốc ưu tiên lô gần hết hạn nhất theo nguyên tắc FEFO. \\
\hline
\end{longtable}

\section{Ma trận truy vết yêu cầu rút gọn}

\begin{longtable}{|p{1.6cm}|p{3.3cm}|p{2.4cm}|p{3.8cm}|p{3cm}|}
\caption{Ma trận truy vết yêu cầu rút gọn}\\
\hline
\textbf{REQ} & \textbf{Use case} & \textbf{Phase} & \textbf{Backend/API} & \textbf{Frontend} \\
\hline
REQ-01 & UC01 Đăng nhập & P1 & POST /api/v1/auth/login & LoginPage \\
\hline
REQ-03 & UC02 Quản lý tài khoản & P1 & /api/v1/users/* & UserManagementPage \\
\hline
REQ-04,05 & UC04-05 Bệnh nhân & P1 & /api/v1/patients & PatientList/Create/Detail \\
\hline
REQ-06,07,08 & UC06-08 Lượt khám & P1 & /api/v1/visits & VisitList/Create \\
\hline
REQ-09-13 & UC09-13 Khám và kê đơn & P1 & /api/v1/examinations & ExaminationPage \\
\hline
REQ-14-16 & UC14-16 Hóa đơn và thanh toán & P1 & /api/v1/invoices, /payments & InvoiceList/Detail \\
\hline
REQ-17-20 & Quy định, danh mục, báo cáo & P1 & regulations, diseases, drugs, reports & Regulation, Catalog, Report pages \\
\hline
REQ-21-30 & UC21-30 mở rộng & P2 & appointments, queue, vitals, services, lab, inventory, pharmacy, organization, audit & P2 feature pages \\
\hline
\end{longtable}

\section{Demo các chức năng trọng tâm}

\begin{longtable}{|p{3.2cm}|p{5.2cm}|p{5.5cm}|}
\caption{Demo chức năng trọng tâm của chương đặc tả yêu cầu}\\
\hline
\textbf{Demo} & \textbf{Nội dung trình bày} & \textbf{Minh chứng cần chèn} \\
\hline
Demo UC tạo lượt khám & Lễ tân chọn bệnh nhân, hệ thống tạo visit và cấp số thứ tự & Use case diagram + screenshot form tạo lượt khám \\
\hline
Demo UC khám bệnh & Bác sĩ mở phiên khám, ghi chẩn đoán, kê đơn & Use case description + screenshot phiếu khám \\
\hline
Demo UC thanh toán & Thu ngân lập hóa đơn, ghi nhận thanh toán & Use case description + screenshot hóa đơn \\
\hline
Demo phân quyền & Mỗi vai trò có quyền và sidebar khác nhau & Bảng RBAC + screenshot sidebar theo role \\
\hline
\end{longtable}

\TwoDemoFigures
{figures/screenshots/visit-create.png}
{Tạo lượt khám}
{figures/screenshots/examination-page.png}
{Lập phiếu khám}
{Minh chứng hai use case lõi trong quy trình khám bệnh}
{fig:req-demo-core}
